% JETSPIN EoM documentation

\section{Equations of Motion\label{sec:Equation-of-Motion}}

The model implemented in JETSPIN is an extension of the Lagrangian
discrete model introduced by Reneker et al. \cite{reneker2000bending}
(in the following text referred as Ren model). The model provides
a compromise of efficiency and accuracy by representing the filament
as a series of $n$ beads (jet beads) at mutual distance $l$ connected
by viscoelastic elements. The length $l$ is typically larger than
the cross-sectional radius of the filament, but smaller than the characteristic
lengths of other observables of interest (e.g. curvature radius).
Each $i-th$ bead has mass $m_{i}$ and charge $q_{i}$ (not necessarily
equal for all the beads). The nanofiber is modeled as a viscoelastic
fluid following the approach developed in Ref. \cite{pontrelli2014electrospinning},
and the stress $\sigma_{i}$ for the element connecting bead $i$
with bead $i+1$ is given by the viscoelastic constitutive equation:

\begin{equation}
\frac{d\sigma_{i}}{dt}=-\frac{1}{t_{0}}\left(\sigma_{i}-\sigma_{HB}\right),\label{eq:huang-garcia}
\end{equation}
where $t_{0}$ is the relaxation time and $\sigma_{HB}$ is the Herschel\textendash Bulkley
stress. Note that the relaxation time is give as

\begin{equation}
t_{0}=\frac{\mu}{G},\label{eq:relaxation-time}
\end{equation}
where $G$ is the elastic modulus and $\mu$ the viscosity of the
fluid jet, while $\sigma_{HB}$ reads

\begin{equation}
\sigma_{HB}=\sigma_{Y}+K_{\upsilon e}\left(\frac{1}{l_{i}}\frac{dl_{i}}{dt}\right)^{n_{\upsilon e}}.\label{eq:Herschel=00003D002013Bulkley}
\end{equation}
In the previous expression, $K_{\upsilon e}$ is the flow consistency
index with dimensions $\text{g}\text{s}^{ne-2}\text{cm}^{-1}$, $\sigma_{Y}$
the yield stress, $n_{\upsilon e}$ the flow behavior index, $t$
the time, and $l_{i}$ the length of the element connecting bead $i$
with bead $i+1$. Note that it is possible to define an effective
viscosity $\mu_{eff}$ as

\begin{equation}
\mu_{eff}=K_{\upsilon e}\left(\frac{1}{l_{i}}\frac{dl_{i}}{dt}\right)^{n_{\upsilon e}-1}.\label{eq:effective-viscosity}
\end{equation}
It is worth stressing that the case $K_{\upsilon e}=\mu$, $n_{\upsilon e}=1$
and $\sigma_{Y}=0$ recovers the Maxwellian fluid model. Keeping the
last case in mind, for convenience sake, we write

\begin{equation}
K_{\upsilon e}=k_{\upsilon e}\mu,\label{eq:rescaled-prefactor}
\end{equation}
with $k_{\upsilon e}$ the flow consistency index re-scaled by $\mu$,
and the viscoelastic constitutive equation is rearranged as

\begin{equation}
\frac{d\sigma_{i}}{dt}=-\frac{G}{\mu}\left[\sigma_{i}-\sigma_{Y}+k_{\upsilon e}\mu\left(\frac{1}{l_{i}}\frac{dl_{i}}{dt}\right)^{n_{\upsilon e}}\right],\label{eq:stress-ode}
\end{equation}
The last constitutive equation allows the modeling of several fluids
as, for example, Bingham fluids ( $\sigma_{Y}>0$ ), shear-thinning
fluids ( $n_{\upsilon e}>1$ ), shear-thickening fluid ( $n_{\upsilon e}<1$
) .

As alternative, it is possible to activate the Kelvin\textendash Voigt
model \cite{barnes1989introduction} (see Tab \ref{tab:inputlist}).
In the last case, Eq \ref{eq:stress-ode} is replaced by

\begin{equation}
\frac{d\sigma_{i}}{dt}=G\frac{1}{l_{i}}\frac{dl_{i}}{dt}+\mu\frac{1}{l_{i}}\frac{d\upsilon_{i}}{dt}.\label{eq:stress-ode-KV}
\end{equation}

Given $a_{i}$ the cross-sectional radius of the filament at the bead
$i$, the viscoelastic force $\vec{\textbf{f}}_{\upsilon e}$ pulling
bead $i$ back to $i-1$ and towards $i+1$ is

\begin{equation}
\vec{\textbf{f}}_{\upsilon e,i}=-\pi a_{i}^{2}\sigma_{i}\cdot\vec{\textbf{t}}_{i}+\pi a_{i+1}^{2}\sigma_{i+1}\cdot\vec{\textbf{t}}_{i+1}\text{,}
\end{equation}
where $\vec{\textbf{t}}_{i}$ is the unit vector pointing bead $i$
from bead $i-1$. The surface tension force $\vec{\textbf{f}}_{st}$
for the $i-th$ bead is given by

\begin{equation}
\vec{\textbf{f}}_{st,i}=k_{i}\cdot\pi\left(\frac{a_{i}+a_{i-1}}{2}\right)^{2}\alpha\cdot\vec{\textbf{c}}_{i}\text{,}
\end{equation}
where $\alpha$ is the surface tension coefficient, $k_{i}$ is the
local curvature, and $\vec{\textbf{c}}_{i}$ is the unit vector pointing
the center of the local curvature from bead $i$. Note the force $\vec{\textbf{f}}_{st}$
is acting to restore the rectilinear shape of the bending part of
the jet.

In the electrospinning experimental configuration an intense external
electric potential $\Phi_{0}$ is applied between the spinneret and
a conducting collector located at distance $h$ from the injection
point. As consequence, each $i-th$ bead endures the electric force

\begin{equation}
\vec{\textbf{f}}_{el,i}=e_{i}\frac{\Phi_{0}}{h}\cdot\vec{\textbf{x}}\text{,}
\end{equation}
where $\vec{\textbf{x}}$ is the unit vector pointing the collector
from the spinneret and assumed along the $x$ axis. Note that in Ren
model the electric field $\Phi_{0}$ is assumed to be static in order
to avoid the computationally expensive integration of Poisson equation,
whereas in reality $\Phi$ is depending on the net charge of the nanofiber
so as to maintain constant the potential at the electrodes. The latter
issue was elegantly addressed by Kowalewski et al. \cite{kowalewski2009modeling},
and its implementation in JETSPIN will be planned. Further, it is
possible to impose a time dependent external electric potential described
by a rectangular waveform with amplitude $\Phi_{0}$ and frequency
given in input file (see Section \ref{sec:Input-file}).

The net Coulomb force $\vec{\textbf{f}}_{c}$ acting on the $i-th$
bead from all the other beads is given by

\begin{equation}
\vec{\textbf{f}}_{c,i}=\sum_{\substack{j=1\\
j\neq i
}
}^{n}\frac{q_{i}q_{j}}{R_{ij}^{2}+a_{j}^{2}}\cdot\vec{\textbf{u}}_{ij}\text{,}
\end{equation}
where $R_{ij}=\left[\left(x_{i}-x_{j}\right)^{2}+\left(y_{i}-y_{j}\right)^{2}+\left(z_{i}-z_{j}\right)^{2}\right]^{1/2}$,
and $\vec{\textbf{u}}_{ij}$ is the unit vector pointing the $i-th$
bead from $j-th$ bead. Note that $j-th$ charge is assumed to induces
a field from its outer shell of the fiber (a charged ring of radius
$a_{j}$) and not from the center line (the particle-like point) as
implemented in the original Ren model. Although the Ren model provides
a reasonable description for the spiral motion of the jet, the last
term $\vec{\textbf{f}}_{c}$ introduces mathematical inconsistencies
due to the discretization of the fiber into point-charges. Different
approaches were developed to overcome this issue which usually imply
strong approximations\cite{feng2002stretching,feng2003stretching,hohman2001electrospinning}.
Other strategies use a less crude approximation by accounting for
the actual electrostatic form factors between two interacting sections
of a charged fiber\cite{kowalewski2005experiments} or involving more
sophisticated methods which exploit the tree-code hierarchical force
calculation algorithm\cite{kowalewski2009modeling}. In JETSPIN it
is possible to activate the dynamic refinement algorithm to recover
a finer discretization of the jet (see Sec \ref{sec:Dynamic-refinement}).

Although usually much smaller than the other driving forces, the body
force due to the gravity is computed in the model by the usual expression

\begin{equation}
\vec{\textbf{f}}_{g,i}=m_{i}g\cdot\vec{\textbf{x}},
\end{equation}
where $g$ is the gravitational acceleration.

As shown in experimental evidences \cite{spinning1991science}, the
air drag force affects the dynamics of the nanofiber. The extended
stochastic model developed in Refs \cite{lauricella2016refinement,Lauricella2015dissipative,lauricella2015nonlinear}
was included in JETSPIN to reproduce the aerodynamic effects. In particular,
the air drag is modeled by adding two force terms: a random term and
a dissipative term. The dissipative air drag term is usually dependent
on the geometry of the jet, which changes in time, and it combines
longitudinal and lateral components. Based on experimental results\cite{spinning1991science},
the longitudinal component of the air drag dissipative force term
acting on bead $i$ is

\begin{equation}
\vec{\textbf{f}}_{air,i}=-0.65\pi a_{i}l_{i}\rho_{a}\left(2a_{i}/\nu_{a}\right)^{-0.81}\upsilon_{tan,i}^{1+0.19}\cdot\vec{\textbf{t}}_{i-1}
\end{equation}
where $\rho_{a}$ and $\nu_{a}$ are the air density and kinematic
viscosity, respectively, $\vec{\textbf{t}}_{i-1}$ is the unit vector
pointing bead $i$ from bead $i-1$, and $\upsilon_{tan,i}=\vec{\boldsymbol{\upsilon}_{i}}\cdot\vec{\textbf{t}}_{i-1}$
is the projection of the vector velocity of bead $i$ along the unit
vector $\vec{\textbf{t}}_{i-1}$. Further, it is possible to write
$\vec{\textbf{f}}_{air,i}$ as

\begin{equation}
\vec{\textbf{f}}_{air,i}=-m_{i}\gamma_{i}\,l_{i}^{0.905}\upsilon_{t,i}^{1+0.19}\cdot\vec{\textbf{t}}_{i-1}
\end{equation}
if we introduced the dissipative coefficient $\gamma_{i}$ equal to

\begin{equation}
\gamma_{i}=0.65\pi^{0.905}\frac{\rho_{a}}{m_{i}}\left(\frac{2}{\nu_{air}}\right)^{-0.81}V_{i}^{0.095},\label{eq:fric_dissip}
\end{equation}
where $V_{i}=\pi a_{i}^{2}\,l_{i}$ is the jet volume represented
by the $i-th$ bead.

Following the expression introduced by A. Yarin\cite{yarin1993free,yarin2014fundamentals},
under a high-speed air drag the lateral component $\vec{\textbf{f}}_{lift,i}$
of the aerodynamic dissipative force related to the flow speed is
given in the linear approximation (for small bending perturbations)
by
\begin{equation}
\vec{\textbf{f}}_{lift,i}=-l_{i}\cdot k_{i}\rho_{a}\upsilon_{i}^{2}\pi\left(\frac{a_{i}+a_{i-1}}{2}\right)^{2}\cdot\vec{\textbf{c}}_{i},
\end{equation}
where $k_{i}$ is still the local curvature, and $\vec{\textbf{c}}_{i}$
is the unit vector pointing the center of the local curvature from
bead $i$. The combined action of such longitudinal and lateral components
provide the dissipative force term acting on the $i-th$ bead

\begin{equation}
\vec{\textbf{f}}_{diss,i}=\vec{\textbf{f}}_{air,i}+\vec{\textbf{f}}_{lift,i}\:.
\end{equation}
Instead, the random force term for the $i-th$ bead has the form

\begin{equation}
\vec{\textbf{f}}_{rand,i}=\sqrt{2m_{i}^{2}D_{\upsilon}}\cdot\vec{\boldsymbol{\eta}_{i}}(t),
\end{equation}
where $D_{\upsilon}$ denotes a generic diffusion coefficient in velocity
space (which is assumed constant and equal for all the beads), and
$\vec{\boldsymbol{\eta}_{i}}$ is a three-dimensional vector, whereof
each component $\eta$ is an independent stochastic process that is
nowhere differentiable with $<\eta\left(t_{1}\right)\eta\left(t_{2}\right)>=\delta\left(\left|t_{2}-t_{1}\right|\right)$,
and $<\eta\left(t\right)>=0$. Note that for the sake of simplicity
we assume $\eta=d\omega(t)/dt$, where $\omega(t)$ is a Wiener process.
Further details are available in Ref \cite{lauricella2016flow}.

The combined action of these forces governs the elongation of the
jet according to the Newton's equation providing a non-linear Langevin-like
stochastic differential equation:

\begin{equation}
m_{i}\frac{d\vec{\boldsymbol{\upsilon}}_{i}}{dt}=\vec{\textbf{f}}_{el,i}+\vec{\textbf{f}}_{c,i}+\vec{\textbf{f}}_{\upsilon e,i}+\vec{\textbf{f}}_{st,i}+\vec{\textbf{f}}_{g,i}+\vec{\textbf{f}}_{diss,i}+\vec{\textbf{f}}_{rand}\:\text{,}\label{eq:force-EOM}
\end{equation}
where $\vec{\boldsymbol{\upsilon}}_{i}$ is the velocity of the $i-th$
bead. The velocity $\vec{\boldsymbol{\upsilon}}_{i}$ satisfies the
kinematic relation:

\begin{equation}
\frac{d\vec{\textbf{r}}_{i}}{dt}=\vec{\boldsymbol{\upsilon}}_{i}\label{eq:pos-EOM}
\end{equation}
where $\vec{\textbf{r}}_{i}$ is the position vector of the $i-th$
bead, $\vec{\textbf{r}}_{i}=x_{i}\vec{\textbf{x}}+y_{i}\vec{\textbf{y}}+z_{i}\vec{\textbf{z}}$.
The three Eqs \ref{eq:stress-ode}, \ref{eq:force-EOM} and \ref{eq:pos-EOM}
form the set of equations of motion (EOM) governing the time evolution
of system. It is worth stressing that Eq \ref{eq:force-EOM} recovers
a deterministic EOM by imposing $\rho_{a}=0$ and $D_{\upsilon}=0$.
